\chapter{第十五讲\quad 刑罚的裁量}
\label{lecture:15}
\noindent{\small\hyperref[ref:lecture:15]{查阅本讲扩展规则}}\par

\section{第一节\quad 量刑情节}
\paragraph{从法定刑档到宣告刑}
设 \(F\) 为已认定的定罪与量刑事实，\(I(F)=[a,b]\) 为由分则和总则确定的本案适用刑档。普通的从轻与从重在 \(I(F)\) 内选择，减轻原则上到 \(a\) 以下；法律规定数个幅度时，在适用档的下一幅度量刑（第61至63条）。无本法规定的减轻情节但案件特殊，须经最高人民法院核准，才有第63条第二款的例外。免除处罚在认定有罪之后适用。符号上分别写 \(\text{应当}(Q)\)、\(\text{可以}(Q)\)、\(\text{不得}(Q)\)，以区分义务、裁量权限和禁止，三者不相互替换。

\paragraph{例题推导}
题设本案适用的有期徒刑档为五年至十年，且唯一从宽依据是法定从轻。由 \(I=[5,10]\) 和第62条推出 \(L\in I\)，故 \(L=4\) 年不符合该规则。若题设另有应当减轻的法定情节，并给出下一档为三年至五年，则先适用第63条，再从下一档取低于本案法定最低刑的期限；四年在候选中，具体数值仍须裁量。一个事实已经用于升档时，再将同一事实无依据地作独立加重，会使评价重复。

\section{第二节\quad 累犯}
\paragraph{一般条件和特殊条件}
令 \(c_0,c_1\) 为前后罪，\(t_0\) 为前罪刑罚执行完毕或赦免日；前罪经假释的，一般累犯的五年从假释期满日起算。以 \(I(c)\) 表示故意罪，\(A(c)\) 表示实施时已满十八周岁，\(P_0\) 表示前罪实际判处有期徒刑以上刑罚，\(P_1\) 表示后罪应判同等级刑罚。第65条一般累犯模型为
\[
 \mathrm{GRec}\iff
 I(c_0)\land I(c_1)\land A(c_0)\land A(c_1)
 \land P_0\land P_1\land T_5(t_0,t_1),
\]
其中 \(T_5\) 以执行完毕或赦免为前提，且后罪发生在其后五年内。这里 \(P_1\) 是本案应判的刑罚，不只看抽象法定最高刑。第66条特殊累犯检验前后罪均属于危害国家安全犯罪、恐怖活动犯罪、黑社会性质的组织犯罪之一，以及前罪刑罚已经执行完毕或经赦免。三类可交叉，无一般累犯的五年与有期徒刑门槛。年龄边界应单独说明：第66条未复写第65条的年龄例外；教材标注的多数观点仍将前后罪犯罪时均已满十八周岁列为条件，属于体系解释口径。

\paragraph{例题推导}
甲成年后因故意犯罪被判有期徒刑，假释后又犯后罪，后罪为成年故意犯罪且应判有期徒刑。若新罪发生于假释考验期内，前罪尚未执行完毕，\(T_5\) 失败，应转查第86、71条；若假释顺利期满两年后才犯，\(T_5\) 成立，全部合取成立，依法从重，并受第74、81条的缓刑、假释禁限。反例：原判缓刑顺利期满意味着原判刑罚不再执行（第76条），不能用其期满日冒充刑罚执行完毕日。

\section{第三节\quad 自首}
\paragraph{先定位哪一项罪行}
对罪行 \(c\) 令 \(V(c)\) 表示自动投案，\(F(c)\) 表示如实供述主要犯罪事实。普通自首入口为 \(V(c)\land F(c)\)。对已受强制措施或正在服刑的法定主体，检验 \(F(c)\)、司法机关尚未掌握本人其他罪行，以及司法解释所规定的余罪范围。是否为不同种罪行通常以罪名区分，选择性罪名或法律、事实密切关联情形还须依司法口径判断（最高法《关于处理自首和立功若干具体问题的意见》）。被动到案后供述已掌握罪行，通常转查坦白。按第67条，自首可以从轻或者减轻，犯罪较轻的还可以免除；坦白可以从轻，因如实供述避免特别严重后果发生的，可以减轻。

\paragraph{例题推导}
甲因盗窃被采取强制措施，主动如实供述机关尚未掌握、与盗窃无关联的另一次故意伤害。就伤害罪，依次得到法定主体、未掌握、另种本人罪行和如实供述，依第67条第二款以自首论；盗窃罪没有自动投案的事实，不随伤害罪获得自首。若仅供述已经掌握的盗窃事实，则“尚未掌握的其他罪行”失败，可能构成的是该盗窃罪的坦白。若记录只说“接受问询”，无从证明自动投案，应说明该事实待查。

\section{第四节\quad 立功}
\paragraph{贡献的类型及证明}
依第68条，立功可由揭发他人犯罪经查证属实、提供侦破其他案件的重要线索等实现；依法协助抓捕也可被认定为立功。令 \(D\) 表示具体揭发他人犯罪且经查实，\(L\) 表示重要线索与侦破有实际联系，\(A\) 表示符合法定认定口径的实际抓捕协助，\(Q\) 表示线索与行为符合司法解释的来源、主体等限制。课堂检查式为
\[
 Q\land(D\lor L\lor A)\ \Longrightarrow\ \text{普通立功候选}.
\]
重大立功另查法定、司法解释条件。普通立功可以从轻或者减轻，重大立功可以减轻或者免除。最高法上述《意见》区分提供自然掌握的同案人基本信息和实质协助，要求对案件侦破、实际抓捕等效果进行核实。

\paragraph{例题推导}
甲到案后按公安机关安排联络乙并将其约至指定地点，公安机关因此抓获乙；协助事实已经查实且无来源排除。约至地点对抓捕有实际作用，成立 \(A\)，结合 \(Q\) 认定立功。若仅按要求提供共同犯罪人姓名，并如实说明自己的共同作案事实，这属于供述范围，不能仅凭姓名推出 \(A\) 或揭发其他犯罪 \(D\)。

\section{第五节\quad 数罪并罚}
\paragraph{第69条的区间与异种刑}
对同种有限刑的各宣告刑 \(l_1,\ldots,l_n\)，令 \(m=\max_i l_i\)，\(s=\sum_i l_i\)。在第69条通常并罚情形，执行刑 \(x\) 满足
\[
 m\le x\le \min\{s,C\},\qquad
 C=\begin{cases}
 3\text{年},&\text{管制},\\
 1\text{年},&\text{拘役},\\
 20\text{年},&\text{有期徒刑且 }s<35\text{年},\\
 25\text{年},&\text{有期徒刑且 }s\ge35\text{年}.
 \end{cases}
\]
定性在前、算区间在后，区间内酌情选择。死刑、无期徒刑走吸收路径；有期徒刑与拘役执行有期徒刑；有期徒刑或拘役与管制并存，先执行前者再执行管制。附加刑仍执行，同种合并、异种分别。法律的“以上”“以下”包括本数（第99条），三十五年恰属二十五年上限支路。若以前案并罚后的刑或余刑参与再并，使机械套式出现下界高于上界，须按题设和适用解释具体处理，不能从空区间随意取值。

\paragraph{第70、71条的时间分流}
在前案刑罚执行完毕前，令前案只判一罪有期徒刑 \(p\) 年、已执行 \(e\) 年，新增罪宣告刑 \(q\) 年。令 \(\mathcal P(a,b)\) 为第69条可选择的执行刑集合。
\[
\begin{array}{l|c|c}
\text{情形}&\text{还须执行}&\text{本案累计实际执行}\\ \hline
\text{宣告前所犯漏罪，第70条}&x-e,\ x\in\mathcal P(p,q)&x\\
\text{宣告后执行完毕前犯新罪，第71条}&y,\ y\in\mathcal P(p-e,q)&e+y
\end{array}
\]
先判断\emph{犯罪发生}与宣告的先后，再判断是否满足适用阶段。“先并后减”中 \(e\) 只减一次，“先减后并”中原刑进入并罚的只是未执行部分。若前案本来就是多罪并罚，第70条并罚基数存在原执行刑与各罪宣告刑等口径讨论；上式特设前案仅一罪。题中兼有漏罪、新罪时，教材按先处理漏罪、扣已执行，再处理新罪的顺序拆算。

\paragraph{例题推导}
原判十四年，已执行六年，新增罪判十二年。若属漏罪，\(m=14\)、\(s=26\)、\(C=20\)，则 \(x\in[14,20]\)；例如取 \(18\)，余 \(18-6=12\) 年。若属新罪，先算 \(14-6=8\)，再有 \(y\in[12,20]\)；若取 \(18\)，这 \(18\) 年是余刑，累计为 \(6+18=24\) 年。单从数字相同不能判断采用哪个公式。反例：新增罪判管制时，不可把“管制一年”与“有期徒刑十四年”混进上述同种刑年数公式。

\section{第六节\quad 缓刑}
\paragraph{资格、强制性、考验期}
第72、74条的入口是被判拘役或三年以下有期徒刑，且非累犯、非犯罪集团首要分子。再合取犯罪情节较轻、有悔罪表现、无再犯危险、对所居住社区无重大不良影响。四项满足后，通常可以宣告；其中未满十八周岁者、怀孕妇女和已满七十五周岁者应当宣告。各年龄事实按适用规则及题意确定。第73条设 \(l\) 为原判刑期、\(p\) 为考验期限：
\[
 p\in
 \begin{cases}
 [\max(l,2\text{个月}),\,1\text{年}],&\text{拘役},\\
 [\max(l,1\text{年}),\,5\text{年}],&\text{有期徒刑}.
 \end{cases}
\]
考验期从判决确定之日起算。附加刑仍执行，可以依法附禁止令。

\paragraph{例题推导}
甲为符合特殊主体条件的怀孕妇女，判有期徒刑十个月，四项实质条件均成立，且无禁止情形。十个月进入有期徒刑支路，满足资格；由特殊主体规则推出应当宣告；下限 \(\max(10\text{个月},1\text{年})=1\text{年}\)。题设若改成判四年，刑种期限门槛失败，怀孕也不能独立触发应当缓刑。考验期满无第77条规定情形，原判刑罚不再执行；期间犯新罪或发现漏罪，撤销缓刑并依第69条合并决定；有严重违规情形的，撤销并执行原判刑罚。

\section*{记忆与闭卷复述}
记“先定档、查再犯、供己首、助案功、漏先并、新先减、缓看四”。拿到题先画宣告、判决确定、开始执行、执行完毕四个点；给从宽事由标注“可以”或“应当”；算并罚写出 \(m,s,C\) 再选择；最后把四个缓刑实质条件逐个对事实。
\begin{enumerate}
 \item \textbf{自测。}三罪各判十五、十、十年有期徒刑，并罚区间是什么？\quad\textbf{核对。}总和三十五年，跨入二十五年上限，最高单刑十五年，故区间为十五至二十五年。
 \item \textbf{自测。}甲被逮捕后才如实供述本案已掌握的罪行，并提供新线索，公安机关据此破获他案。两个从宽事由各是什么？\quad\textbf{核对。}前者按坦白审查，后者按立功审查；评价单位和证明对象不同。
 \item \textbf{自测。}甲满足缓刑四项条件且审判时怀孕，但属于累犯。应否宣告？\quad\textbf{核对。}第74条排除，特殊主体的强制性规则须在资格成立后适用。
\end{enumerate}
